\section{Pressure-Change Formulation}
To force the surrogate to emulate actual fluid dynamics rather than exploiting temporal autocorrelation, we reformulate the regression problem to target the one-step pressure increment:
\begin{equation}
\label{eq:delta_p_def}
\Delta P_i(t) = P_i(t+1) - P_i(t) = f(X_{\text{local}, i}(t))
\end{equation}
This formulation aligns the regression objective with the spatial-temporal derivative terms in the discretized flow equation (Eq.~\ref{eq:fvm}). Under this target, the Persistence baseline (which predicts $\Delta P = 0$) scores an $R^2 = -0.0198$ with an RMSE of 4.22~bar, exposing the true difficulty of the task (Table~\ref{tab:baseline_results}). The linear OLS model achieves an $R^2 = 0.0007$ (RMSE = 4.18~bar), which is statistically indistinguishable from the Persistence baseline. The MLP model reaches only an $R^2 = 0.3524$ and an RMSE of 3.36~bar. These results reveal the genuine difficulty of emulating pressure changes using cell-local inputs.

\subsection{Local Feature Insufficiency}
The near-zero performance improvement of the linear OLS model over the Persistence baseline highlights the inadequacy of local cell features. To investigate this, a leave-one-out feature ablation study was conducted on the training set, measuring the change in test set RMSE when each input feature was sequentially removed. Quantitatively, the resulting $\Delta \text{RMSE}$ was bounded within $\pm 0.0007$~bar for all cell-wise features, including absolute permeability, cell porosity, current pressure $P_i(t)$, and the dynamic well control parameters.

This result provides strong evidence that cell-local features are insufficient for emulating transient flow. A physical explanation is provided by Eq.~\ref{eq:fvm}: the pressure change at any cell $i$ is driven by the transmissibility-weighted pressure differences with its adjacent neighbors and by localized well source terms. Both spatial pressure gradients and local well indices are absent from the baseline cell-wise feature set. A purely local model has no access to the spatial gradients driving Darcy flow, and without temporal history, it cannot characterize the transient decay that follows changes in well control.
